\documentclass[12pt]{article}

% Idioma y codificación
\usepackage[spanish, es-tabla]{babel}       %es-tabla para que se titule "Tabla"
\usepackage[utf8]{inputenc}

% Márgenes
\usepackage[a4paper,top=3cm,bottom=2.5cm,left=3cm,right=3cm]{geometry}

% Comentarios de bloque
\usepackage{verbatim}

% Paquetes de links
\usepackage[hidelinks]{hyperref}    % Permite enlaces
\usepackage{url}                    % redirecciona a la web

% Más opciones para enumeraciones
\usepackage{enumitem}

% Personalizar la portada
\usepackage{titling}


% Paquetes de tablas
\usepackage{multirow}


%------------------------------------------------------------------------

%Paquetes de figuras
\usepackage{caption}
\usepackage{subcaption} % Figuras al lado de otras
\usepackage{float}      % Poner figuras en el sitio indicado H.


% Paquetes de imágenes
\usepackage{graphicx}       % Paquete para añadir imágenes
\usepackage{transparent}    % Para manejar la opacidad de las figuras

% Paquete para usar colores
\usepackage[dvipsnames]{xcolor}
\usepackage{pagecolor}      % Para cambiar el color de la página

% Habilita tamaños de fuente mayores
\usepackage{fix-cm}


%------------------------------------------------------------------------

% Paquetes de matemáticas
\usepackage{mathtools, amsfonts, amssymb, mathrsfs}
\usepackage[makeroom]{cancel}     % Simplificar tachando
\usepackage{polynom}    % Divisiones y Ruffini
\usepackage{units} % Para poner fracciones diagonales con \nicefrac

\usepackage{pgfplots}   %Representar funciones
\pgfplotsset{compat=1.18}  % Versión 1.18

\usepackage{tikz-cd}    % Para usar diagramas de composiciones
\usetikzlibrary{calc}   % Para usar cálculo de coordenadas en tikz

% Paquete para usar el símbolo de grados
\usepackage{gensymb}

%Definición de teoremas, etc.
\usepackage{amsthm}
%\swapnumbers   % Intercambia la posición del texto y de la numeración

\theoremstyle{plain}

\makeatletter
\@ifclassloaded{article}{
  \newtheorem{teo}{Teorema}[section]
}{
  \newtheorem{teo}{Teorema}[chapter]  % Se resetea en cada chapter
}
\makeatother

\newtheorem{coro}{Corolario}[teo]           % Se resetea en cada teorema
\newtheorem{prop}[teo]{Proposición}         % Usa el mismo contador que teorema
\newtheorem{lema}[teo]{Lema}                % Usa el mismo contador que teorema

\theoremstyle{remark}
\newtheorem*{observacion}{Observación}

\theoremstyle{definition}

\makeatletter
\@ifclassloaded{article}{
  \newtheorem{definicion}{Definición} [section]     % Se resetea en cada chapter
}{
  \newtheorem{definicion}{Definición} [chapter]     % Se resetea en cada chapter
}
\makeatother

\newtheorem*{notacion}{Notación}
\newtheorem*{ejemplo}{Ejemplo}
\newtheorem*{ejercicio*}{Ejercicio}             % No numerado
\newtheorem{ejercicio}{Ejercicio} [section]     % Se resetea en cada section


% Modificar el formato de la numeración del teorema "ejercicio"
\renewcommand{\theejercicio}{%
  \ifnum\value{section}=0 % Si no se ha iniciado ninguna sección
    \arabic{ejercicio}% Solo mostrar el número de ejercicio
  \else
    \thesection.\arabic{ejercicio}% Mostrar número de sección y número de ejercicio
  \fi
}


% \renewcommand\qedsymbol{$\blacksquare$}         % Cambiar símbolo QED
%------------------------------------------------------------------------

% Paquetes para encabezados
\usepackage{fancyhdr}
\pagestyle{fancy}
\fancyhf{}

\newcommand{\helv}{ % Modificación tamaño de letra
\fontfamily{}\fontsize{12}{12}\selectfont}
\setlength{\headheight}{15pt} % Amplía el tamaño del índice


%\usepackage{lastpage}   % Referenciar última pag   \pageref{LastPage}
\fancyfoot[C]{\thepage}

%------------------------------------------------------------------------

% Conseguir que no ponga "Capítulo 1". Sino solo "1."
\makeatletter
\@ifclassloaded{book}{
  \renewcommand{\chaptermark}[1]{\markboth{\thechapter.\ #1}{}} % En el encabezado
    
  \renewcommand{\@makechapterhead}[1]{%
  \vspace*{50\p@}%
  {\parindent \z@ \raggedright \normalfont
    \ifnum \c@secnumdepth >\m@ne
      \huge\bfseries \thechapter.\hspace{1em}\ignorespaces
    \fi
    \interlinepenalty\@M
    \Huge \bfseries #1\par\nobreak
    \vskip 40\p@
  }}
}
\makeatother

%------------------------------------------------------------------------
% Paquetes de cógido
\usepackage{minted}
\renewcommand\listingscaption{Código fuente}

\usepackage{fancyvrb}
% Personaliza el tamaño de los números de línea
\renewcommand{\theFancyVerbLine}{\small\arabic{FancyVerbLine}}

% Estilo para C++
\newminted{cpp}{
    frame=lines,
    framesep=2mm,
    baselinestretch=1.2,
    linenos,
    escapeinside=||
}


\usepackage{listings} % Para incluir código desde un archivo

\renewcommand\lstlistingname{Código Fuente}
\renewcommand\lstlistlistingname{Índice de Códigos Fuente}

% Definir colores
\definecolor{vscodepurple}{rgb}{0.5,0,0.5}
\definecolor{vscodeblue}{rgb}{0,0,0.8}
\definecolor{vscodegreen}{rgb}{0,0.5,0}
\definecolor{vscodegray}{rgb}{0.5,0.5,0.5}
\definecolor{vscodebackground}{rgb}{0.97,0.97,0.97}
\definecolor{vscodelightgray}{rgb}{0.9,0.9,0.9}

% Configuración para el estilo de C similar a VSCode
\lstdefinestyle{vscode_C}{
  backgroundcolor=\color{vscodebackground},
  commentstyle=\color{vscodegreen},
  keywordstyle=\color{vscodeblue},
  numberstyle=\tiny\color{vscodegray},
  stringstyle=\color{vscodepurple},
  basicstyle=\scriptsize\ttfamily,
  breakatwhitespace=false,
  breaklines=true,
  captionpos=b,
  keepspaces=true,
  numbers=left,
  numbersep=5pt,
  showspaces=false,
  showstringspaces=false,
  showtabs=false,
  tabsize=2,
  frame=tb,
  framerule=0pt,
  aboveskip=10pt,
  belowskip=10pt,
  xleftmargin=10pt,
  xrightmargin=10pt,
  framexleftmargin=10pt,
  framexrightmargin=10pt,
  framesep=0pt,
  rulecolor=\color{vscodelightgray},
  backgroundcolor=\color{vscodebackground},
}

%------------------------------------------------------------------------

% Comandos definidos
\newcommand{\bb}[1]{\mathbb{#1}}
\newcommand{\cc}[1]{\mathcal{#1}}

% I prefer the slanted \leq
\let\oldleq\leq % save them in case they're every wanted
\let\oldgeq\geq
\renewcommand{\leq}{\leqslant}
\renewcommand{\geq}{\geqslant}

% Si y solo si
\newcommand{\sii}{\iff}

% Letras griegas
\newcommand{\eps}{\epsilon}
\newcommand{\veps}{\varepsilon}
\newcommand{\lm}{\lambda}

\newcommand{\ol}{\overline}
\newcommand{\ul}{\underline}
\newcommand{\wt}{\widetilde}
\newcommand{\wh}{\widehat}

\let\oldvec\vec
\renewcommand{\vec}{\overrightarrow}

% Derivadas parciales
\newcommand{\del}[2]{\frac{\partial #1}{\partial #2}}
\newcommand{\Del}[3]{\frac{\partial^{#1} #2}{\partial^{#1} #3}}
\newcommand{\deld}[2]{\dfrac{\partial #1}{\partial #2}}
\newcommand{\Deld}[3]{\dfrac{\partial^{#1} #2}{\partial^{#1} #3}}


\newcommand{\AstIg}{\stackrel{(\ast)}{=}}
\newcommand{\Hop}{\stackrel{L'H\hat{o}pital}{=}}

\newcommand{\red}[1]{{\color{red}#1}} % Para integrales, destacar los cambios.

% Método de integración
\newcommand{\MetInt}[2]{
    \left[\begin{array}{c}
        #1 \\ #2
    \end{array}\right]
}

% Declarar aplicaciones
% 1. Nombre aplicación
% 2. Dominio
% 3. Codominio
% 4. Variable
% 5. Imagen de la variable
\newcommand{\Func}[5]{
    \begin{equation*}
        \begin{array}{rrll}
            #1:& #2 & \longrightarrow & #3\\
               & #4 & \longmapsto & #5
        \end{array}
    \end{equation*}
}

%------------------------------------------------------------------------
\setlength{\parindent}{0pt}

\setlist[itemize, 1]{label=\textbullet \textbf{)}}  % Cambiar las viñetas a círculos
\usepackage{xkeyval} % Para el paso de argumentos

\usepackage{graphicx}

% Definir la carpeta de las imágenes
\graphicspath{{../_assets}{../../_assets}}

% Definir el comando \portada
\makeatletter
\define@key{portada}{titulo}{\def\titulo{#1}}
\define@key{portada}{subtitulo}{\def\subtitulo{#1}}
\define@key{portada}{autor}{\def\autor{#1}}
\define@key{portada}{año}{\def\año{#1}}

\newcommand*{\portada}[1][]{%
  % Encabezado
  \fancyhead[L]{\helv \titulo}
  \fancyhead[R]{\helv \nouppercase{\leftmark}}

    % Definimos las claves y sus valores por defecto
  \setkeys{portada}{%
    titulo=Sin Título,%
    subtitulo=Sin Subtítulo,%
    autor=Autor Desconocido,%
    año=Sin Año, #1}%
    \begin{titlepage}
        \centering
        {\includegraphics[width=0.2\textwidth]{Logo-UGR-Black.png}\par}
        \vspace{1cm}
        {\bfseries\LARGE Universidad de Granada \par}
        \vspace{1cm}
        {\scshape\Large Doble Grado en Ingeniería Informática y Matemáticas \par}
        \vspace{3cm}
        {\scshape\Huge \titulo \par}
        \vspace{3cm}
        {\itshape\Large \subtitulo \par}
        \vfill
        {\Large Autor: \par}
        {\Large \autor \par}
        \vfill
        {\Large \año \par}
    \end{titlepage}%
}

\begin{document}
    \portada[%
        titulo=Variable Compleja I,
        subtitulo=Resumen,
        autor=Jesús Muñoz Velasco,
        año=Curso 2024-2025]


    % \section{Test de Weierstrass}

    % Sea $\sum\limits_{n\geq 0} f_n$ una serie de funciones complejas definidas en un conjunto $A\subset \bb{C}$, y sea $B\subset A$. Supongamos que, para cada $n\in \bb{N}\cup \{0\}$, existe una constante $M_n\in \bb{R}$ tal que:
    % \begin{gather*}
    %     |f_n(z)| \leq M_n \hspace{1cm} \forall z \in B
    % \end{gather*}
    % Si la serie de números reales $\sum\limits_{n\geq 0} M_n$ es convergente, entonces la serie de funciones $\sum\limits_{n\geq 0} f_n$ converge absoluta y uniformemente en $B$.

    % \section{Teorema de convergencia de Weierstrass}

    % Sea $\Omega$ un abierto no vacío de $\bb{C}$ y para cada $n\in \bb{N}$, sea $f_n\in \cc{H}(\Omega)$. Supongamos que la sucesión $\{f_n\}$ converge uniformemente en cada subconjunto compacto de $\Omega$ a una función $f:\Omega \to \bb{C}$, en particular:
    % \begin{gather*}
    %     f(z) = \lim_{n\to \infty} f_n(z) \hspace{1cm} \forall z \in \Omega
    % \end{gather*}
    % Entonces $f\in \cc{H}(\Omega)$ y para cada $k\in \bb{N}$ se tiene que la sucesión $\{f_n^{(k)}\}$ de las k-ésimas derivadas, converge a la derivada k-ésima $f^{(k)}$, uniformemente en cada subconjunto compacto de $\Omega$, en  particular:
    % \begin{gather*}
    %     f^{(k)}(z) = \lim_{n\to \infty} f_n^{(k)}(z) \hspace{1cm} \forall z \in \Omega,\ \ \forall k \in \bb{N}
    % \end{gather*}

    % \section{Holomorfía de funciones dependientes de un parámetro}

    \section{Tema 3}

    \begin{teo}
        Sea $A$ un subconjunto no vacío de $\bb{C}$ y $f\in \cc{F}(A)$. Como $A\subset \bb{R}^2$, podemos considerar las funciones $u,v:A \to \bb{R}$ definidas, para todo $(x,y)\in A$ por
        \begin{gather*}
            u(x,y) = Re\ f(x+iy) \hspace{1cm} y \hspace{1cm} v(x,y) = Im\ f(x+iy)
        \end{gather*}
        Para $z_0=(x_0,y_0)\in A^\circ$, las siguientes afirmaciones son equivalentes:
        \begin{enumerate}
            \item La función $f$ es derivable en el punto $z_0$.
            \item Las funciones $u$ y $v$ son diferenciables en el punto $(x_0, y_0)$ verificando que 
            \begin{gather*}
                \frac{\partial u}{\partial x} (x_0, y_0) = \frac{\partial v}{\partial y} (x_0, y_0) \hspace{1cm} y \hspace{1cm} \frac{\partial u}{\partial y} (x_0, y_0) = -\frac{\partial v}{\partial x} (x_0, y_0)
            \end{gather*}
            Estas ecuaciones reciben el nombre de \textbf{ecuaciones de Cauchy-Riemann}.
        \end{enumerate}
        Caso de que se cumplan 1. y 2., se tiene:
        \begin{gather*}
            f'(z_0) = \frac{\partial u}{\partial x} (x_0, y_0) + i \frac{\partial v}{\partial x} (x_0, y_0)
        \end{gather*}
    \end{teo}

    \section{Tema 4}

    \begin{teo}[\textbf{Test de Weierstrass}]
        Sea $\sum\limits_{n\geq 0} f_n$ una serie de funciones complejas definidas en un conjunto $A\subset \bb{C}$, y sea $B\subset A$. Supongamos que, para cada $n\in \bb{N}\cup \{0\}$, existe una constante $M_n\in \bb{R}$ tal que:
        \begin{gather*}
            |f_n(z)| \leq M_n \hspace{1cm} \forall z \in B
        \end{gather*}
        Si la serie de números reales $\sum\limits_{n\geq 0} M_n$ es convergente, entonces la serie de funciones $\sum\limits_{n\geq 0} f_n$ converge absoluta y uniformemente en $B$.
    \end{teo}

    \begin{lema}[\textbf{Lema de Abel}] Dado $\rho\in \bb{R}^+$, supongamos que la sucesión $\{|\alpha_n|\rho^n\}$ está mayorada. Entonces la serie de potencias $\sum\limits_{n\geq 0} \alpha_n ( z-a)^n$ converge absolutamente en el disco abierto $D(a, \rho)$ y uniformemente en cada compacto $K$ que esté contenido en dicho disco.
    \end{lema}

    \begin{prop}[\textbf{Fórmula de Cauchy-Hadamard}]
        Sea $R$ el radio de convergencia de la serie $\sum\limits_{n\geq0}\alpha_nz^n$
        \begin{enumerate}
            \item Si la suceción $\{\sqrt[n]{|\alpha_n|}\}$ no está mayorada, entonces $R=0$.
            \item Si $\{\sqrt[n]{|\alpha_n|}\} \to 0$, entonces $R=\infty$.
            \item En otro caso se tiene: $R = \dfrac{1}{\limsup \{\sqrt[n]{|\alpha_n|}\}}$
        \end{enumerate}
    \end{prop}

    \begin{coro}
        Supongamos que $\alpha_n\in \bb{C}^*$ para todo $n\in \bb{N}$ y sea $R$ el radio de convergencia de la serie de potencias $\sum\limits_{n\geq 0} \alpha_n z^n$.
        \begin{enumerate}
            \item Si $\{\alpha_{n+1} / \alpha_n\} \to \infty$, entonces $R=0$.
            \item Si $\{\alpha_{n+1} / \alpha_n\} \to 0$, entonces $R=\infty$.
            \item Si $\{\alpha_{n+1} / \alpha_n\} \to \lambda \in \bb{R}^+$, entonces $R=1/\lambda$.
        \end{enumerate}
    \end{coro}

    \section{Tema 6}

    \begin{teo}[\textbf{Caracterización de la existencia de primitiva}] Sea $\Omega$ un abierto no vacío de $\bb{C}$ y $f\in \cc{C}(\Omega)$. Las siguientes afirmaciones son equivalentes:
    \begin{enumerate}
        \item Existe $F\in \cc{H}(\Omega)$ tal que $F'(z) = f(z)$ para todo $z\in \Omega$.
        \item Para todo camino cerrado $\gamma$ en $\Omega$ se tiene que $\displaystyle\int_\gamma f(z)dz = 0$.
    \end{enumerate}
        
    \end{teo}

    \section{Tema 7}

    \begin{teo}[\textbf{Teorema local de Cauchy}]
        Si $\Omega$ es un dominio estrellado, entonces toda función admite una primitiva en $\Omega$, es decir, existe $F\in \cc{H}(\Omega)$ tal que $F'(z) = f(z)$ para todo $z\in \Omega$. Equivalentemente se tiene
        \begin{gather*}
            \int_\gamma f(z) dz = 0
        \end{gather*}
        Para toda función $f\in \cc{H}(\Omega)$ y todo camino cerrado $\gamma$ en $\Omega$.
    \end{teo}

    \begin{prop}[\textbf{Fórmula de Cauchy}]
        Sean $\Omega$ un abierto de $\bb{C}$ y $f\in \cc{H}(\Omega)$. Dado $a\in \Omega$, sea $r\in \bb{R}^+$ tal que $\overline{D}(a,r)\subset \Omega$. Se tiene entonces:
        \begin{gather*}
            f(z) = \frac{1}{2\pi i } \int_{C(a,r)} \frac{f(w)}{w-z} dw \hspace{1cm} \forall z \in D(a,r)
        \end{gather*}
    \end{prop}

    \section{Tema 8}
    \begin{teo}[\textbf{Desarrollo en serie de Taylor}]
        Si $\Omega$ es un abierto no vacío de $\bb{C}$ y $f\in \cc{H}(\Omega)$, entonces $f$ es analítica en $\Omega$ y, en particular, $f$ es indefinidamente derivable en $\Omega$. Además:
        \begin{enumerate}
            \item Si $\Omega = \bb{C}$, para todo $a\in \bb{C}$, la serie $\sum\limits_{n\geq 0} \dfrac{f^{(n)}(a)}{n!}(z-a)^n$ tiene radio de convergencia infinito y se verifica que:
            \begin{gather*}
                f(z) = \sum\limits_{n=0}^\infty \frac{f^{(n)}(a)}{n!} (z-a)^n \hspace{1cm} \forall z \in \bb{C}
            \end{gather*}
            \item Si $\Omega \neq \bb{C}$ y para cada $a\in \Omega$ tomamos $R_a = d(a,\bb{C}\setminus \Omega)$, la serie $\sum\limits_{n\geq 0} \dfrac{f^{(n)}(a)}{n!}(z-a)^n$ tiene radio de convergencia mayor o igual que $R_a$ y se verifica que 
            \begin{gather*}
                f(z) = \sum\limits_{n=0}^\infty \frac{f^{(n)}(a)}{n!} (z-a)^n \hspace{1cm} \forall z \in D(a, R_a)
            \end{gather*}
        \end{enumerate}
    \end{teo}

    \begin{prop}[\textbf{Teorema de Cauchy para las derivadas}] Sean $\Omega$ un abierto de $\bb{C}$ y $f\in \cc{H}(\Omega)$. Dado $a\in \Omega$, sea $r\in \bb{R}^+$ tal que $\overline{D}(a,r)\subset \Omega$. Se tiene entonces:
        \begin{gather*}
            f^{(k)} (z) = \frac{k!}{2\pi i } \int_{C(a,r)} \frac{f(w)}{(w-z)^{k+1}} dw \hspace{1cm} \forall z \in D(a,r),\ \forall k \in \bb{N}\cup \{0\}
        \end{gather*}
    \end{prop}

    \begin{teo}[\textbf{Teorema de extensión de Riemann}]
        Sean $\Omega$ un abierto de $\bb{C}$, $z_0 \in \Omega$ y $f\in \cc{H}(\Omega\setminus \{z_0\})$. Las siguientes afirmaciones son equivalentes:
        \begin{enumerate}
            \item Existe $g\in \cc{H}(\Omega)$ tal que $g(z) = f(z)$ para todo $z\in \Omega \setminus \{z_0\}$.
            \item $f$ tiene límite en el punto $z_0$.
            \item Existen $\delta,M>0$ tales que $|f(z)|\leq M$ para todo $z \in \Omega$ que verifique $0 < |z - z_0| < \delta$.
            \item $\lim\limits_{z \to z_0} (z - z_0) f(z) = 0$.
        \end{enumerate}
    \end{teo}

    \section{Tema 9}

    \begin{prop}[\textbf{Desigualdades de Cauchy}] Sean $\Omega$ un abierto de $\bb{C}$, $f\in \cc{H}(\Omega)$ y $a\in \Omega$. Dado $r\in \bb{R}^+$ tal que $\overline{D}(a,r)\subset \Omega$, sea $M(f,a,r) = \max \{|f(z)|: z\in C(a,r)^*\}$. Se tiene entonces:
    \begin{gather*}
        \frac{|f^{(n)}(a)|}{n!} \leq \frac{M(f,a,r)}{r^n} \hspace{1cm} \forall n \in \bb{N}\cup \{0\}
    \end{gather*}       
    \end{prop}

    \begin{teo}[\textbf{Teorema de Liouville}] Toda función entera y acotada es constante. De hecho, la imagen de cualquier función entera no constante es un conjunto denso en $\bb{C}$, es decir, para $f\in \cc{H}(\bb{C})$ tal que $\exists M\in \bb{R}^+$ de forma que $|f(z)| \leq M$ $\forall z \in \bb{C}$, entonces se tiene que $\overline{Im(f)} = \bb{C}$.
    \end{teo}

    \begin{teo}[\textbf{Teorema fundamental de Álgebra}] El cuerpo $\bb{C}$ es algebraicamente cerrado, es decir, si $P$ es un polinomio con coeficientes complejos, no constante, existe $z\in \bb{C}$ tal que $P(z) = 0$.
    \end{teo}

    \begin{prop}[\textbf{Principio de identidad para funciones holomorfas}] Sea $\Omega$ un dominio y $f,g\in \cc{H}(\Omega)$. Si $A$ es un subconjunto de $\Omega$ tal que $f(a) = g(a)$ para todo $a\in A$, y $A' \cap \Omega \neq \emptyset$, entonces $f$ y $g$ son idénticas: $f(z)=g(z)$ para todo $z\in \Omega$.
    \end{prop}

    \section{Tema 10}

    \begin{teo}[\textbf{Teorema de convergencia de Weierstrass}] Sea $\Omega$ un abierto no vacío de $\bb{C}$ y, para cada $n\in \bb{N}$, sea $f_n \in \cc{H}(\Omega)$. Supongamos que la sucesión $\{f_n\}$ converge uniformemente en cada subconjunto compacto de $\Omega$ a una función $f:\Omega \to \bb{C}$, en particular:
    \begin{gather*}
        f(z) = \lim_{n\to \infty} f_n(z) \hspace{1cm} \forall z \in \Omega
    \end{gather*}
    Entonces $f\in \cc{H}(\Omega)$ y, para cada $k\in \bb{N}$, se tiene que la sucesión $\{f^{(k)}_n\}$ de las k-ésimas derivadas, converge a la derivada k-ésima $f^{(k)}$, uniformemente en cada subconjunto compacto de $\Omega$, en particular:
    \begin{gather*}
        f^{(k)} (z) = \lim_{n\to \infty} f_n^{(k)} (z) \hspace{1cm} \forall z \in \Omega,\ \forall k \in \bb{N}
    \end{gather*}
    Este resultado también se puede usar para series
    \end{teo}

    \begin{teo}[\textbf{Holomorfía de la integral dependiente de un parámetro}]
        Sea $\gamma$ u ncamino, $\Omega$ un abierto del plano y $\Phi:\gamma^* \times \Omega \to \bb{C}$ una función continua. Supongamos que, para cada $w\in \gamma^*$, la función $\Phi_w:\Omega \to \bb{C}$ definida por $\Phi_w(z) = \Phi(w,z)$ para todo $z\in \Omega$, es holomorfa en $\Omega$. Entonces, definiendo
        \begin{gather*}
            f(z) = \int_\gamma \Phi(w,z) dw \hspace{1cm} \forall z \in \Omega
        \end{gather*}
        se obtiene una función holomorfa: $f \in \cc{H}(\Omega)$. Además, para cada $k\in \bb{N}$ y cada $z\in \Omega$, la función $w \mapsto \Phi_w ^{(k)}(z)$, de $\gamma^*$ en $\bb{C}$, es continua y se verifica que
        \begin{gather*}
            f^{(k)} (z) = \int_\gamma \Phi_w^{(k)} dw = \int_\gamma \frac{\partial^k \Phi}{\partial z^k}(w,z) \hspace{1cm} \forall z \in \Omega, \forall k \in \bb{N}
        \end{gather*}
    \end{teo}

    \section{Tema 11}

    \begin{prop}[\textbf{Propiedad de la media}]
        Sea $\Gamma$ un abierto de $\bb{C}$ y $f\in \cc{H}(\bb{C})$. Para $a\in \Omega$ y $r\in \bb{R}^+$ tales que $\overline{D}(a,r)\subset \Omega$, se tiene:
        \begin{gather*}
            f(a) = \frac{1}{2\pi} \int_{-\pi}^{\pi} f(a+re^{it}) dt
        \end{gather*}
    \end{prop}

    \begin{teo}[\textbf{Principio del módulo máximo}] Sea $\Omega$ un dominio y $f\in \cc{H}(\Omega)$. Supongamos que $|f|$ tiene un máximo relativo en un punto $a\in \Omega$, es decir, existe $\delta > 0$ tal que $D(a, \delta)\subset \Omega$ y $|f(z)| \leq |f(a)|$ para todo $z\in D(a, \delta)$. Entonces $f$ es constante.
    \end{teo}

    \begin{teo}[\textbf{Principio del módulo mínimo}] Sea $\Omega$ un dominio y $f\in \cc{H}(\Omega)$. Supongamos que $|f|$ tiene un mínimo relativo en un punto $a\in \Omega$ es decir, existe $\delta > 0$ tal que $D(a, \delta)\subset \Omega$ y $|f(z)| \geq |f(a)|$ para todo $z\in D(a, \delta)$. Entonces, o bien $f(a)=0$, o bien $f$ es constante.
    \end{teo}

    \begin{teo}[\textbf{Teorema de la función inversa global}] Sea $U$ un dominio y $f\in \cc{H}(U)$ una función inyectiva. Entonces $V=f(U)$ es un dominio y $f^{-1}\in \cc{H}(V)$ con 
    \begin{gather*}
        (f^{-1})'(f(z)) = \frac{1}{f'(z)} \hspace{1cm} \forall z \in U
    \end{gather*}
    \end{teo}

    \section{Tema 12}

    \begin{teo}[\textbf{Forma general del teorema de Cauchy y la fórmula integral de Cauchy}] Sea $\Omega$ un abierto del plano y $\Gamma$ un ciclo en $\Omega$, nul-homólogo con respecto a $\Omega$. Para toda función $f\in \cc{H}(\Omega)$ se tiene:
    \begin{enumerate}
        \item $\displaystyle Ind_\Gamma (z)f(z) = \frac{1}{2\pi i} \int_\Gamma \frac{f(w)}{w-z} dw$ \hspace{1cm} $\forall z \in \Omega \setminus \Gamma^*$
        \item $\displaystyle \int_\Gamma f(w) dw = 0$
    \end{enumerate}
    \end{teo}

    
    
\end{document}